\section{The coefficient action and its fixed geometry}
\label{sec:coefficients}

The coefficient ring and its fixed deformation geometry supply two
inputs to the topological differential calculation.  The ring model
gives the coefficient cohomology, while the fixed geometry constrains
equivariant maps between its scalar and exterior parts.
We first determine the actual action: the model in
Theorem~\ref{coef:model} below is an isomorphism of completed graded rings
with their actual \(C_9\)-actions.  It is the height-six, prime-three case
of the model proposed in \cite[Conjecture~0.3]{BR}; that conjecture is not
an input to the proof.
Related geometric descriptions at height \(p-1\) are given by
Ray \cite[Theorem~C and Section~10.5]{Ray}; at the prime three
that height is two.  The height-six model and its fixed deformation
geometry are established here from the formal-group action.

Retain the Honda formal group, the chosen subgroup
\(G=\langle g\rangle=C_9\), and \(H=\langle\sigma\rangle\), where
\(\sigma=g^3\), from Section~\ref{sec:detector}.  Put
\[
 k=\F_{3^6},\qquad
 R=W(k)[[u_1,\ldots,u_5]],\qquad
 \m=(3,u_1,\ldots,u_5),\qquad E_*=R[u^{\pm1}],
 \quad |u|=-2.
\]
All completions of graded rings in this section are taken degree by
degree.  The coefficient line \(E_{-2}\) is the Lie line of the universal
formal group: its dual \(E_2\) is the cotangent line, identified by
restriction from \(\mathbf{CP}^{\infty}\) to \(\mathbf{CP}^1\).
We use the right-unit convention for the strict coordinate-change
coefficients \(t_i\).  Thus \(t_1\) represents the usual integral
one-cocycle of internal degree four.

Write
\[
 \mathcal C=C_{\calO_D^\times}(H),\qquad
 \calO_D=W(k)\langle S\rangle/(S^6=3,\ Sa=a^\varphi S).
\]
The action of \(\mathcal C\) considered here fixes \(W(k)\).  The
centralizer algebra of \(\Q_3(\zeta_3)\) is a division algebra of degree
three over that field.  It contains its unramified cubic extension.
Skolem--Noether, followed by adjustment of the valuation of a conjugator
by a normalizing power of \(S\), allows us to arrange that
\[
 T=\F_{27}^{\times}\subseteq\mathcal C
\]
consists of Teichm\"uller units.  We index its characters modulo \(26\),
with the residual character on the Lie line assigned weight one.
Since \(26\) is invertible in \(W(k)\), character projections may be
applied to lifts in all the complete modules below.

\subsection{The tangent action}

To obtain the cyclic presentation, we will construct a period
\(x_0\in E_{-2}\) with \(\Tr_H(x_0)=0\) and show that its
\(G\)-orbit generates the completed coefficient ring.  The tangent
calculation gives the required independent coordinates.  The following
complete-local calculation then lifts the residual period to one with
exactly zero trace.

For an element \(e\) acting on coefficients, let \(D_e=e-1\).  Choose
Araki generators temporarily, so that
\([3](x)=\sum_i^F v_i x^{3^i}\), and normalize the period by
\(v_6=u^{1-3^6}\).  This normalization is possible by Hensel's lemma.
The elements \(u_i=v_i u^{3^i-1}\), \(1\le i\le5\), are regular
parameters.  Set
\[
 P_0=[3],\qquad P_i=[u_i]\quad(1\le i\le5)
\]
in \(\m/\m^2\).

\begin{lemma}[Tangent action]\label{coef:tangent}
For \(1\le i\le5\), the operator \(D_g\) sends \(P_i\) into the span of
\(P_0,\ldots,P_{i-1}\), with a nonzero coefficient of \(P_{i-1}\).
The operator \(D_\sigma\) lowers the index by at least three, and its
coefficient at \(P_{i-3}\), for \(i\ge3\), is
\(a^{3^{i-3}}\), where
\[
 a\,a^{27}=-1.
\]
Moreover, \(D_g u/u\) has a nonzero \(P_5\)-term and
\(D_\sigma u/u\) has leading term \(a^{27}P_3\).
One may choose a weight-one period \(u\) and linear coordinates
\[
 y_0=[D_\sigma u/u],\qquad y_1=P_4,\qquad y_2=P_5,\qquad
 l_i=D_\sigma y_i
\]
such that
\[
 l_0=-P_0,\qquad
 \operatorname{gr}_{\m}E_*
   =k[u^{\pm1},y_0,l_0,y_1,l_1,y_2,l_2],
 \qquad
 \sigma y_i=y_i+l_i,\quad \sigma l_i=l_i,\quad \sigma u=u.
\]
The weights of \(y_0,y_1,y_2\) are \(0,2,8\), respectively.
\end{lemma}

\begin{proof}
For \(e=g,\sigma\), comparison of the two 3-series under the strict
coordinate change gives, through index six,
\begin{equation}\label{coef:right-unit-linear}
 ev_i\equiv v_i+\sum_{l<i}v_l\,t_{i-l}(e)^{3^l}
 \pmod{\m^2}.
\end{equation}
Here and below an expression involving a graded coefficient is read
after multiplication by the appropriate power of the period.
To justify the truncation, in
\(f\circ[3]_{eF}=[3]_F\circ f\) all coefficients \(v_l,ev_l\) for
\(l<6\) belong to \(\m\).  Terms containing two such coefficients
vanish modulo \(\m^2\).  Inside a term already multiplied by \(v_l\)
one may work in characteristic three.  The special-fibre Honda law is
additive below total degree \(3^6\), and terms involving the
index-six summand together with a positive nonlinear power exceed the
degree under consideration.  These observations leave precisely
\eqref{coef:right-unit-linear}.

The \(S\)-valuations of \(g-1\) and \(\sigma-1\) are one and three.
Consequently their first residual nonlinear strict coefficients have
indices one and three, respectively.  The identity
\((\sigma-1)^2=-3\sigma\) in the endomorphism algebra implies
\(a\,a^{27}=-1\) for the latter coefficient.  Equation
\eqref{coef:right-unit-linear} gives the asserted descending terms.
At index six the equation \(v_6=u^{-728}\), together with
\(-728\equiv1\pmod3\), gives the period formulas.

For Teichm\"uller elements the lower strict coefficients vanish on the
special fibre.  It follows that \(P_i\) has weight \(3^i-1\).
Projecting a period to weight one preserves its residual generator and
the leading terms just calculated.  In particular
\[
 D_\sigma y_0=a^{27}aP_0=-P_0.
\]
The six displayed coordinates are independent: the three \(y_i\)
have leading indices \(3,4,5\), and the three \(l_i\) have leading
indices \(0,1,2\).  Their asserted action follows either from this
triangular calculation or from
\(D_\sigma^2=0\) on \(\m/\m^2\).
\end{proof}

\subsection{A complete-local coefficient calculation}

For this subsection only, use \(\mathcal E_r\) for the spectral sequence
of the \(\m\)-adic filtration on the coefficient Tate complex.  This
auxiliary spectral sequence is distinct from the topological
homotopy-fixed-point spectral sequence.  On \(\mathcal E_r\), the
differential increases the \(\m\)-filtration by \(r\), increases
cohomological degree by one, and preserves internal degree.

The complete cyclic Tate complex in a fixed internal degree alternates
\(D_\sigma\) and \(\Tr_H=1+\sigma+\sigma^2\).  Let \(b\) denote its
periodicity class of cohomological degree two and internal degree zero.
Set
\[
 N_i=y_i^3-y_i l_i^2.
\]
The letters \(N_i\) initially denote elements of the associated graded;
their multiplicative orbit norms will provide lifts.

\begin{lemma}[The first coefficient page]\label{coef:graded-tate}
There is an isomorphism of trigraded algebras
\begin{equation}\label{coef:first-page}
 \mathcal E_1
 =k[b^{\pm1},u^{\pm1},N_0,N_1,N_2]
       \otimes\Lambda(e,z_0,z_1,z_2).
\end{equation}
Here \(e,z_i\) have cohomological degree one and internal degree zero.
Their cyclic cocycle values at \(\sigma\) are \(1,y_i\), respectively.
Their \(\m\)-filtrations are \(0,1\), and \(N_i\) has filtration three.
The classes \(b^{\pm1}\), \(N_i\), and the class denoted \(u^3\) lift
to cycles of the complete coefficient complex.
\end{lemma}

\begin{proof}
For one pair \(k[y,l]\) with \(\sigma y=y+l\), the orbit norm is
\(N=y^3-yl^2\), and the ring is free over \(k[N,l]\) on \(1,y,y^2\).
After removing powers of \(N\), its homogeneous part of degree \(d\ge2\)
has basis
\[
 l^d,\qquad yl^{d-1},\qquad y^2l^{d-2}.
\]
The action on this three-dimensional space is a single unipotent block
of length three; it is a free \(k[H]\)-module.  Degree zero is trivial,
and degree one is the two-dimensional indecomposable module, a
suspension of the trivial module in the stable category.

Tensor the three decompositions.  Terms having a free factor have zero
Tate cohomology.  The stable suspension class of each degree-one factor
is represented by its value \(y_i\), and cup product with that class
identifies the corresponding shifted Tate groups.  This gives
\eqref{coef:first-page} as a module, and also identifies products on
distinct factors.  Odd anticommutation and the vanishing of odd
squares give the displayed exterior algebra.

The elements \(N_i\) lift by taking multiplicative orbit norms of
lifts of \(y_i\); the norm of a period has initial term \(u^3\).
The periodicity operator \(b\) and its inverse act on the actual cyclic
complex.  These are therefore coefficient cycles.
\end{proof}

\begin{proposition}[Integral obstruction calculation]
\label{coef:integral-lift}
After compatible nonzero scalar normalizations, the auxiliary
coefficient spectral sequence has
\[
 d_1(u)=uz_0,\qquad
 d_1(e)=ez_0,\qquad
 d_1(z_i)=z_i z_0\ (i=1,2),\qquad d_1(z_0)=0.
\]
Consequently
\begin{equation}\label{coef:second-page}
 \mathcal E_2
 =k[b^{\pm1},u^{\pm3},N_0,N_1,N_2]
       \otimes\Lambda(eu,z_1u,z_2u,z_0).
\end{equation}
Its only nonzero differential on these generators is
\[
 d_2(z_0)=c\,bN_0,\qquad c\in k^\times.
\]
There are no further differentials.  In particular the leading
cohomology class \(eu\) is represented by an actual element
\(x_0\in E_{-2}\) satisfying
\[
 x_0/u\equiv1\pmod{\m},\qquad
 \Tr_H(x_0)=0.
\]
The element \(x_0\) may be chosen of \(T\)-weight one.
\end{proposition}

\begin{proof}
The definition of \(y_0\) gives \(d_1(u)=uz_0\).  The linear part of
\(\Tr_H(u)/u\) is \(P_0+l_0=0\), so \(d_1(eu)=0\).  In the cyclic
cup-product convention this is equivalently \(d_1(e)=ez_0\).

We need the same trace assertion to one further order.  The trace of
the chosen period now lies in \(\m^2E_{-2}\).  It is invariant, so its
leading term defines an even Tate class in filtration two.  The only
classes there with the required internal degree are the shifts of
\(uz_i z_j\), \(i<j\), by powers of \(b\).  None has \(T\)-weight one:
the added weights are \(2,8,10\).  The leading trace class is therefore
zero.  Subtracting the trace of an element of \(\m^2E_{-2}\) corrects
\(u\) so that
\[
 \Tr_H(u)\in\m^3E_{-2}.
\]
Since \(g\) commutes with \(H\), the elements \(D_g u,D_g^2u,D_\sigma u\)
satisfy the same trace condition.  By
Lemma~\ref{coef:tangent} their first terms span
\(uz_2,uz_1,uz_0\), in triangular order, modulo \(l_i\)-terms.
Those \(l_i\)-terms are boundaries in the odd cyclic group.
Thus \(d_1(z_1u)=d_1(z_2u)=0\).  The equation \(d_1^2(u)=0\)
also gives \(d_1(z_0)=0\).
Leibniz yields the asserted formulas on \(e,z_1,z_2\).
Taking cohomology, a monomial survives precisely when its power of
\(u\) is congruent modulo three to its number of factors among
\(e,z_1,z_2\); this gives \eqref{coef:second-page}.

The three generators \(eu,z_1u,z_2u\) all have period exponent one
modulo three.  An even target at that exponent must contain exactly
one of these generators and also \(z_0\).  For a \(d_2\) from \(z_1u\)
or \(z_2u\) this is impossible in filtration three; for a \(d_2\) from
\(eu\), the filtration-two candidates have the wrong weight.
For \(z_0\), a target has period exponent zero and filtration three.
The norms \(N_1,N_2\) have weights \(6,24\), whereas \(N_0\) has
weight zero.  The remaining possible four-exterior-factor target,
with its period corrected by \(u^{-3}\), has weight \(10\).
Thus the only possible differential is \(d_2z_0=c\,bN_0\).

To determine \(c\), let \(L=\Q_3(\zeta_3)\) and take a height-one
Lubin--Tate formal module for the unramified cubic extension of \(L\).
Its underlying formal group has height six.  After a residue-field
extension, choose a framing identifying the action of \(\zeta_3\)
with the prescribed \(H\)-action.  Skolem--Noether gives such a
framing rationally; it can be made integral because the centralizer
division algebra contains an element of \(S\)-valuation one.
This construction uses formal-module deformation theory and the
given stabilizer embedding, not the coefficient model being proved.

On this fixed lift the ungraded coefficients have trivial \(H\)-action
and the Lie period has character \(\zeta_3\), or its inverse if the
opposite action convention is chosen.  Put \(\pi=\zeta_3-1\).
Filter the target cyclic complex by powers of \(\pi\).
The class \(z_0\) maps to a nonzero multiple of \(\pi e\).
Since \(3\) is a unit times \(\pi^2\), the norm differential on this
internal-degree-zero class sends \(\pi e\) to a nonzero multiple of
\(\pi^3b\).  This is its \(d_2\), and \(N_0\) maps to \(\pi^3\)
up to a unit.  Naturality proves \(c\ne0\).

The resulting page is
\begin{equation}\label{coef:final-graded}
 k[b^{\pm1},u^{\pm3},N_1,N_2]
       \otimes\Lambda(eu,z_1u,z_2u).
\end{equation}
Its three odd generators have no even target of period exponent one
modulo three.  Its even generators already lift to coefficient
cycles.  Hence all subsequent differentials vanish.

We finally justify passage to the complete complex.  In each fixed
internal and cohomological degree its terms are compact: their
\(\m\)-adic quotients are finite because \(k\) is finite.
Survival to every page means that every prescribed finite accuracy
admits a lift with the given leading cohomology class.  A boundary
correction in the leading quotient permits us to fix a representative
of that class.  The sets of such lifts whose differentials vanish to
accuracy \(N\) are nonempty nested closed subsets of a compact module.
Their intersection supplies an exact cycle.  Images of the
differential, and images of deeper filtration pieces, are closed for
the same reason.  Thus the final page really is the associated
graded of the coefficient Tate cohomology; no inverse-limit
extension is being suppressed.

Apply this argument to \(eu\).  An odd cyclic cocycle is precisely an
element in the kernel of \(\Tr_H\), so it gives the claimed \(x_0\).
Projection to the weight-one character commutes with \(\Tr_H\),
preserves the residual period, and gives the last assertion.
\end{proof}

\begin{theorem}[The exact cyclic coefficient model]\label{coef:model}
Set \(x_i=g^i x_0\), with indices read modulo nine.  Then
\begin{equation}\label{coef:ring-model}
 \left(
 \frac{W(k)[x_0,\ldots,x_8,(x_0\cdots x_8)^{-1}]}
 {(x_i+x_{i+3}+x_{i+6}:0\le i<3)}
 \right)^{\wedge}_{(3,x_i/x_0-1)}
 \xrightarrow{\ \cong\ } E_*
\end{equation}
is an isomorphism of \(W(k)\)-linear graded \(G\)-rings, with
\(|x_i|=-2\) and \(g(x_i)=x_{i+1}\).
\end{theorem}

\begin{proof}
The actual trace equation of Proposition~\ref{coef:integral-lift}
and its first two \(g\)-translates give the three relations.
All \(x_i/x_0\) reduce to one, so the product being inverted is a unit.

Eliminate \(x_6,x_7,x_8\), and use \(x_0\) and the five ratios
\(x_i/x_0\), \(1\le i\le5\), as coordinates.  The degree-zero part
of the source is a power-series ring in five variables over \(W(k)\);
the remaining variable is an invertible period.
The tangent map is invertible.  Indeed the first five positive
differences \(D_g^j x_0/x_0\) have leading terms in
\(P_{6-j}\), respectively, with nonzero coefficients.  Any
first-order correction from \(u\) to the weight-one \(x_0\), after
division by \(u\), has weight zero and lies in the span of \(P_0,P_3\).
Its positive \(D_g\)-differences only change lower terms in this
triangular calculation.  Together with \(3\), the five displayed
differences therefore give all six cotangent directions.
The passage from these differences to the five ratio coordinates
is an invertible triangular binomial change of variables.
Formal inversion over \(W(k)\) proves that the degree-zero map is an
isomorphism, and adjoining the period proves
\eqref{coef:ring-model}.  Its equivariance holds by the definition
\(x_i=g^i x_0\).
\end{proof}

\subsection{Coefficient cohomology and its scalar classes}

Return now to the bidegree convention \((s,t)\), with group cohomological
degree first and internal homotopy degree second.  The following is the
coefficient algebra that forms the topological \(E_2\)-page after
Euler localization.

\begin{proposition}\label{coef:h-cohomology}
Put \(q_i=\Nm_H(x_i)\), \(q=q_0\).  There are norm parameters \(n_1,n_2\)
such that, with \(K=k[[n_1,n_2]]\),
\begin{equation}\label{coef:h-tate-ring}
 \widehat H^*(H;E_*)=
 K[q^{\pm1},b^{\pm1}]\otimes_K\Lambda_K(V),
 \qquad |q|=(0,-6),\quad |b|=(2,0),\quad |V|=(1,4).
\end{equation}
The module \(V\) is free of rank three, with cyclic-cocycle basis
\[
 a_i=x_i/q_i,\qquad 0\le i<3.
\]
The group \(g\) cyclically permutes this basis semilinearly over \(K\).
Equivalently, \(q_i/q-1\), \(i=1,2\), are regular parameters of \(K\).
For \(\mathfrak n=(n_1,n_2)\), the weights on
\(\mathfrak n/\mathfrak n^2\) are \(6,-2\), and \(V/\mathfrak n V\)
is a single three-term Jordan block for \(g\).
\end{proposition}

\begin{proof}
Lift the norm classes \(N_1,N_2\) in \eqref{coef:final-graded}, choosing
weight \(2,8\) lifts before taking norms.  Their norms have weights
\(6,24\).  The trace-zero \(x_i\) give odd cocycles, and division by
\(q_i\) places them in internal degree four.  The first two
\(g\)-differences yield the \(z_2u,z_1u\) terms of
\eqref{coef:final-graded}; together with \(x_0\), they give the three
odd generators.  One may first use a common denominator \(q\) to
verify this associated-graded assertion.

Tate cohomology is killed by three, so it is a \(k\)-algebra.
All the power series in the positive-filtration norm parameters
converge in its complete filtration.  The proposed generators
account for the entire associated graded.  Their only proposed
relations are the exterior relations, which hold in cohomology by
graded commutativity and invertibility of two.  The complete
associated-graded comparison therefore proves
\eqref{coef:h-tate-ring}.

For the alternative parameters, the identity
\[
 \Nm_H(a+b)-\Nm_H(a)-\Nm_H(b)\in\im(\Tr_H)
\]
holds in every commutative \(H\)-ring: the mixed terms in the triple
product are the two complete cyclic orbits.  Apply it to
\(a=1\), \(b=x_i/x_0-1\).
Modulo additive traces, \(q_i/q-1\) is the norm of that ratio
difference.  The first two ratio differences span the \(y_1,y_2\)
directions modulo the \(y_0,l_i\) directions, by the triangular
tangent calculation; the norms of the latter directions vanish on
\eqref{coef:final-graded}.  Thus these two norm functions are another
regular parameter system.

Finally \(q_{i+3}=q_i\), and \(x_{i+3}-x_i=D_\sigma x_i\) is an odd
cyclic boundary.  Hence \(g(a_2)=a_0\) in cohomology.
Modulo \(\mathfrak n\), semilinearity disappears and the cyclic
permutation of three basis elements is a single Jordan block in
characteristic three.
\end{proof}

\begin{lemma}[Scalar cycles]\label{coef:scalar-cycles}
On every nonzero topological Tate page,
\[
 A=qb^{-2}
\]
is a permanent-cycle unit, and every element of \(K\) is a cycle.
\end{lemma}

\begin{proof}
The representation-graded multiplicative norm of \(x_0\) is an actual
class of degree minus twice the real regular representation of \(H\).
Its symbol is \(q u_\lambda^2\).  Division by the Euler square,
whose symbol is \(b^2u_\lambda^2\), gives \(A\).  Both numerator and
Euler class have actual filtered lifts.  Their quotient is a unit
on each nonzero localized page and has no outgoing differential.

The ratios
\[
 q_i/q=\Nm_H(x_i/x_0)
\]
are norms of elements of \(\pi_0E\), so they have degree-zero
homotopy lifts.  The norm parameters used in
Proposition~\ref{coef:h-cohomology} have the same property.
Constants are cycles: a Teichm\"uller constant satisfies
\(c^{729}=c\), and Leibniz kills its differential on a
characteristic-three page.  For an arbitrary series use the finite
decomposition
\[
 f=\sum_{0\le i,j<3}f_{ij}^{\,3}n_1^i n_2^j,\qquad f_{ij}\in K.
\]
Leibniz kills each cube and each parameter.  This proves the assertion
without a continuity assumption on a spectral-sequence derivation.
\end{proof}

\subsection{The fixed formal-module deformation}

We now study the centralizer action on the scalar and exterior modules
in \eqref{coef:h-tate-ring}. To constrain equivariant maps between them,
we evaluate their classes on deformations fixed by \(H\), and then on a
one-dimensional locus in that fixed space. This requires identifying
the fixed functor on all Artinian tests, including nonreduced ones.
We begin with the finite-height rigidity underlying
Lubin--Tate deformation theory; compare
\cite[Proposition~2.4]{LubinTate}.

\begin{lemma}[Rigidity of endomorphism lifts]\label{coef:rigidity}
Let \(B\) be an Artinian local ring with residue characteristic three,
and let \(F_1,F_2\) be deformations over \(B\) of formal groups of
finite height.  A homomorphism \(F_1\to F_2\) is determined by its
reduction.  The same assertion holds over a complete local base by
passage to its Artinian quotients.
\end{lemma}

\begin{proof}
The formal-group difference of two homomorphisms with the same
reduction is a homomorphism \(f\) whose coefficients lie in the maximal
ideal \(\mathfrak b\) of \(B\).  Suppose inductively that they lie in
\(\mathfrak b^i\), \(i\ge1\).  In the equality
\[
 [3]_{F_2}\circ f=f\circ[3]_{F_1},
\]
the left side is zero modulo \(\mathfrak b^{i+1}\): its linear term
is \(3f\), and every nonlinear term has at least two factors from
\(\mathfrak b^i\).  On the right, only substitution by the
special-fibre 3-series matters.  That series has positive finite
order, so substitution by it is injective on
\((\mathfrak b^i/\mathfrak b^{i+1})[[x]]\).
Thus \(f\) vanishes modulo \(\mathfrak b^{i+1}\).
Induction proves \(f=0\).
\end{proof}

Write
\[
 \calO=\Z_3[\zeta],\qquad \zeta=\zeta_3,\qquad \pi=\zeta-1.
\]
A \emph{strict formal \(\calO\)-module} means a formal group with
\(\calO\)-action whose induced action on its Lie line is the
structural scalar action of \(\calO\).  For formal-module deformation
theory we use the Lubin--Tate--Drinfeld theorem: the framed deformation
of a one-dimensional formal \(\calO\)-module of relative height \(d\)
is a formally smooth space of relative dimension \(d-1\) over the
completed unramified coefficient extension of \(\calO\).
This is the strict formal-module deformation problem; see
\cite{LubinTate,Drinfeld} and \cite[Theorem~8.4]{RapoportZink}.
The normalization of the uniformizer coefficients will be used
separately below, with its own reference and coordinate comparison.

For a coefficient line we shall write \(w^i\) for the \(i\)-th tensor
power of the Lie line, retaining its \(\mathcal C\)-equivariant
structure.  Thus a section written \(f w^i\) has scalar coefficient
\(f\) only after a choice of the displayed period.
Put \(P_n=Kq^n\); after a quotient of \(K\), the same notation denotes
the corresponding scalar line over that quotient.

\begin{theorem}[Fixed geometry and evaluation]\label{coef:fixed-geometry}
The eigenvalue condition
\[
 \sigma x_i=\zeta x_i
\]
defines a framed strict formal \(\calO\)-module deformation space of
relative height three.  Its ungraded base and Lie period are
\[
 D_0=W(k)[\zeta][[r_1,r_2]],\qquad
 r_i=x_i/x_0-1\ (i=1,2),\qquad w=x_0.
\]
The full ungraded \(H\)-fixed base is equivalently
\begin{equation}\label{coef:nonreduced-fixed}
 W(k)[[s,r_1,r_2]]/(s^2+3s+3),\qquad 1+s=\sigma x_0/x_0.
\end{equation}
In particular its characteristic-three fibre retains the relation
\(s^2=0\).

Set \(\bar D=D_0/(\pi)\).  There is a
\(\mathcal C\)-invariant Hasse section \(h_1\) in the line \(w^{-2}\),
given by the first strict coefficient of \(\sigma\).
On its zero divisor there is a \(\mathcal C\)-invariant Hasse section
\(h\) in \(w^{-8}\), given by the second strict coefficient.
The ring
\[
 D'=\bar D/(h_1)
\]
is a formal disc with parameter \(hw^8\); here ideals generated by
line sections are read after a period trivialization.

Evaluation in cyclic cohomology and reduction modulo \(\pi\) identify
\[
 K\cong\bar D^{\,3},\qquad q\longmapsto w^3.
\]
The evaluated \(V\) is the \(\bar D^3\)-span of
\(1,x_1/x_0,x_2/x_0\) in the line \(w^{-2}\).
There is a \(\mathcal C\)-fixed primitive element
\(\theta\in V\), the restriction of the class of \(t_1\), which
evaluates to \(h_1\) up to a nonzero scalar.

One may choose a \(T\)-basis \(v,\theta,y\) of \(V\) with weights
\(-2,0,6\), with \(v=x_0/q\).  Modulo \(\mathfrak n\),
\begin{equation}\label{coef:jordan}
 g\theta=\theta,\qquad
 gy=y+c_1\theta,\qquad
 gv=v+c_2y+c_3\theta,\qquad c_1c_2\ne0.
\end{equation}
Using evaluation to identify cube sections with their unique
preimages, put
\[
 J=h_1^3\in Kq^{-2},\qquad F=q^2J,\qquad
 S_0=K/(F)\cong(D')^3,\qquad
 W=(V/K\theta)\otimes_K S_0.
\]
Then \(F\) is a regular parameter of \(K\), and there is an injective
\(\mathcal C\)-equivariant evaluation
\[
 W\longrightarrow w^{-2}\quad\text{over }D'.
\]
After trivializing the line, the evaluations of \(v,y\) have
orders zero and one, respectively.
\end{theorem}

\begin{proof}
We first compute the fixed equations, without passing to reduced
points.  In the model of Theorem~\ref{coef:model}, put
\(a_i=x_i/x_0\), so \(a_0=1\), and put \(r=a_3\).
Fixation of the ungraded base by \(\sigma\) says
\[
 a_{i+3}=r a_i.
\]
It follows that \(a_4=ra_1\), \(a_5=ra_2\), and \(a_6=r^2\).
The first trace relation also gives \(a_6=-1-r\).  Thus
\(r^2+r+1=0\).  Conversely, this equation and the displayed formulas
impose all the base-fixation equations, including the remaining two
trace relations.  Writing \(r=1+s\) gives
\eqref{coef:nonreduced-fixed}.  Fixing its structural root \(r=\zeta\)
gives precisely \(D_0\).  The action of \(\mathcal C\) preserves this
locus because it centralizes \(\sigma\).

On an Artinian test ring, a fixed deformation carries the lift of
the framed automorphism \(\sigma\).  This lift is unique by
Lemma~\ref{coef:rigidity}.  Its endomorphism
\(1+\sigma+\sigma^2\) reduces to zero and is therefore zero.
It consequently extends the \(\Z_3\)-action to an
\(\calO\)-action.  The eigenvalue condition imposes exactly the
structural action on the Lie line.  Conversely, such a strict
formal-module deformation supplies this lifted automorphism and
the eigenvalue condition.  This proves the claimed identification
of deformation functors on all Artinian tests.
Since \([\calO:\Z_3]=2\), its underlying height six is relative
height three.

On \(\bar D\) the structural uniformizer vanishes.  The first
nonlinear coefficient of the uniformizer endomorphism is
\(t_1(\sigma)\), a section of \(w^{1-3}=w^{-2}\).
After this coefficient is zero, its next leading coefficient is
\(t_2(\sigma)\), a section of \(w^{1-9}=w^{-8}\).
Subtracting the identity from \(\sigma\) gives these statements
successively modulo the previous leading coefficients.
Leading coefficients of an endomorphism transform in the indicated
lines under coordinate change.  They are therefore intrinsic
\(\mathcal C\)-invariant sections, denoted \(h_1,h\).
Here is the parameter comparison explicitly.  After the faithfully
flat completed residue extension \(k\subseteq\bar k\), the universal
formal-module deformation admits parameters \(U_1,U_2\) and a formal
group coordinate \(z\) satisfying
\begin{equation}\label{coef:height-coordinates}
 [\pi](z)\equiv U_i z^{3^i}
 \pmod{(\pi,U_1,\ldots,U_{i-1},z^{3^i+1})},
 \qquad i=1,2;
\end{equation}
this is \cite[arXiv version, Section~1, Equation~(1.1)]{StrauchTorsion}, applied to
the ramified local field \(\Q_3(\zeta)\).
If \(z'=cz+O(z^2)\), with \(c\) a unit, conjugation sends a first
nonzero coefficient \(\lambda z^N\) to \(c^{1-N}\lambda\).
Thus \(h_1w^2\) is a unit times \(U_1\), and, on its zero divisor,
\(hw^8\) is a unit times \(U_2\).
Nonvanishing of these successive cotangent classes descends along
\(k\subseteq\bar k\).  The formal inverse function theorem therefore
makes \(h_1w^2\) a regular parameter of \(\bar D\), and identifies
its zero curve with \(k[[hw^8]]\).  This proves the disc assertion
over the original finite residue field; no descent of the chosen
normal-form coordinate is required.

At the fixed deformation,
\(\Nm_H(x_i)=x_i^3\), since \(1+\zeta+\zeta^2=0\) and
\(\zeta^{0+1+2}=1\).  Hence
\[
 q_i/q\longmapsto (x_i/x_0)^3=(1+r_i)^3.
\]
The parameter description of
Proposition~\ref{coef:h-cohomology}, and perfection of \(k\), identify
\(K\) with \(k[[r_1^3,r_2^3]]=\bar D^3\).
Using the common denominator \(q\), the values of the three odd
generators are \(w^{-2}\), \((1+r_1)w^{-2}\), and
\((1+r_2)w^{-2}\).  The alternative denominators \(q_i\) multiply
these by units from the cube subring, so give the same span.

The class of \(t_1\) is the restriction of the integral
Adams--Novikov one-cocycle and is fixed by the centralizer.
Its value is \(h_1\), up to the nonzero scalar determined by the
cyclic generator convention.  Its evaluated first jet is nonzero,
so its class \(\theta\) is primitive in \(V\).
The residual weights of \(V\), from
\eqref{coef:final-graded} after division by the norm period, are
\(-2,0,6\).  Character projection gives a basis \(v,\theta,y\)
of these weights.  The tangent chain in
Lemma~\ref{coef:tangent} puts \(\theta\) at the fixed end and gives
\eqref{coef:jordan}, including nonvanishing of \(c_1,c_2\).

Under the identification \(K=\bar D^3\),
\(F=q^2h_1^3\) is \((h_1w^2)^3\); it is a parameter.
This also gives \(S_0\cong(D')^3\).
The values and first derivatives of \(1,r_1,r_2\) span first jets.
The class \(\theta\) is transverse to the curve, while \(v\) has
nonzero constant value.  The remaining class \(y\) has nonzero
derivative along the curve.  Equivalently its scalar first jet has
weight \(8\), whereas the transverse parameter has weight \(2\).
Thus, on the curve, the evaluated \(v,y\) have orders zero and one.

For completeness, these orders prove the asserted injectivity over
the smaller ring.  Write \(\tau=hw^8\), so
\(D'=k[[\tau]]\) and \(S_0=k[[\tau^3]]\).
A nonzero cube-series multiple of the value of \(v\) has order
\(0\pmod3\), and one of the value of \(y\) has order \(1\pmod3\).
They cannot cancel.  Evaluation of \(W\) is therefore injective.
All constructions used to define it are functorial for the
centralizer action, proving equivariance.
\end{proof}

\subsection{Centralizer rigidity on the curve}

\begin{lemma}\label{coef:infinite-action}
Every open subgroup of \(\mathcal C\) has infinite image on \(D'\).
Its invariant Laurent functions are constants.  These assertions
remain valid after a finite extension of the residue field, with
the group acting trivially on the extended constants.
\end{lemma}

\begin{proof}
Suppose an open subgroup \(U\) fixed the curve.
Its elements would lift to automorphisms of the formal
\(\calO\)-module over \(D'\), with their given special fibres.
The relative height at the generic point is two:
the first height parameter is zero and the second is nonzero.
After extending the fraction field to an algebraic closure, the
centralizer of \(\Q_3(\zeta)\) in the rational endomorphism algebra
therefore has dimension four over \(\Q_3(\zeta)\).

At the closed point the corresponding centralizer is a division
algebra of degree three, hence dimension nine over \(\Q_3(\zeta)\).
An open subgroup of its units spans this algebra: differences of
elements in a sufficiently small neighborhood of one contain a
nonzero multiple of its full integral lattice.
Choose nine members of \(U\) linearly independent over
\(\Q_3(\zeta)\).
Their generic lifts would satisfy a nontrivial linear relation
over that field of operators.  Clear its denominators using the
structural \(\calO\)-action.  The resulting sum is an integral
formal-module endomorphism; if it vanishes generically, all its
power-series coefficients vanish over the domain \(D'\).
Specialization gives the same forbidden relation at the closed
point.  This proves that no open subgroup fixes the curve.
A finite image would have an open kernel, since distinct
automorphisms of a formal disc are distinguished on finite jets.
Thus every open subgroup has infinite image.

If a Laurent function is nonconstant, subtract its constant term
when necessary and invert it when it has a pole.  Its leading order
shows that the Laurent-series field of the disc is finite over the
one-variable Laurent-series field it generates.  A group of field
automorphisms fixing that function is consequently finite.
It cannot contain the image of an open subgroup of \(\mathcal C\).
The same argument and the same dimension contradiction apply after
a finite extension of constants.
\end{proof}

\begin{lemma}[Invariant meromorphic sections]\label{coef:hasse-powers}
A nonzero invariant meromorphic section of \(w^i\) is a constant
multiple of \(h^d\), where \(8d=-i\) and \(d\in\Z\).
If it is integral, then \(d\ge0\).
There is no nonzero invariant meromorphic section of \(w^i\chi\)
for any nontrivial constant sign character
\(\chi:\mathcal C\to\{\pm1\}\).
\end{lemma}

\begin{proof}
Let \(z\) be such a section.  Since \(h\) is invariant in \(w^{-8}\),
the product \(z^8h^i\) is an invariant Laurent function.
It is a nonzero constant by Lemma~\ref{coef:infinite-action}.
In a period trivialization, \(h\) has order one.
Valuation therefore gives \(8\ord(z)+i=0\).
Put \(d=\ord(z)\); then \(i=-8d\), and \(z/h^d\) has eighth
power constant.  After a finite extension of constants it is itself
constant.  This proves the first assertion and its integral variant.

For a sign twist the same eighth-power calculation applies.
The resulting constant \(z/h^d\) would have to transform by the
nontrivial character \(\chi\), although the group fixes constants.
It must therefore vanish, a contradiction.
\end{proof}

\subsection{The determinant character}

The next calculation is over the reduced ring \(\bar D\), not over
the nonreduced characteristic-three fibre in
\eqref{coef:nonreduced-fixed}.  This distinction is needed when
dividing the determinant of a crystalline operator.
The use of Kodaira--Spencer theory to express an equivariant
canonical line in terms of a Lie line and a determinant character
has a close antecedent in the Lubin--Tate calculation of
Hopkins--Gross \cite[Section~1.2 and Corollary~3]{HopkinsGross}.
Here the strict \(\calO\)-action and the quotient by the invariant
Hasse direction require the following calculation on the curve.

\begin{lemma}[The determinant line]\label{coef:determinant}
After a finite extension of constants there is an equivariant
isomorphism
\begin{equation}\label{coef:determinant-formula}
 \det(D'\otimes_{S_0}W)\cong w^{-9}\chi,
\end{equation}
where \(\chi:\mathcal C\to\{\pm1\}\) is a sign character whose
restriction to \(T\) has weight \(13\), and is therefore nontrivial.
\end{lemma}

\begin{proof}
Write \(A_0=\bar D\), and let \(\mathcal X/A_0\) be its universal
strict formal \(\calO\)-module.
Let \(\mathcal M\) denote the covariant de Rham Dieudonn\'e module
of its underlying \(3\)-divisible group.  Our convention is
\(\mathcal M=\mathbb D(\mathcal X^D)\), where \(\mathbb D\) is the
contravariant crystal and \(\mathcal X^D\) is the Cartier dual.
The crystal and its duality are provided by
\cite[Section~2.2, Theorem~2, and Section~5.2, Theorem~5(i)]{BerthelotMessing}.
It is free of rank six
over \(A_0\).  The uniformizer endomorphism acts on \(\mathcal M\)
by an operator \(\epsilon\) with
\(\epsilon^2=-3(1+\epsilon)=0\).
We first show that
\[
 \mathcal M\ \text{is free of rank three over }\
 B_0=A_0[\epsilon]/(\epsilon^2).
\]

At the closed point, the integral Dieudonn\'e lattice is
torsion-free over \(W(k)\otimes_{\Z_3}\calO=W(k)[\zeta]\).
Indeed multiplication by three is injective, and the relation
\(\pi^2=-3(1+\pi)\) implies that multiplication by \(\pi\) is
injective as well.  This coefficient ring is a discrete valuation
ring, so the lattice is free of rank three over it.
Reduction modulo three gives a rank-three free module over
\(k[\epsilon]/(\epsilon^2)\).
Lift such a basis to \(\mathcal M\).  The resulting
\(B_0\)-linear map \(B_0^3\to\mathcal M\) is a map of free
\(A_0\)-modules of rank six which is an isomorphism at the closed
point.  Its determinant is a unit, proving the assertion.

Set \(M_0=\mathcal M/\epsilon\mathcal M\), and let \(L\) be the Lie
line, denoted \(w\).  The Hodge quotient factors through \(M_0\)
because the structural uniformizer acts as zero on \(L\).
Write
\begin{equation}\label{coef:hodge-quotient}
 0\longrightarrow N\longrightarrow M_0\longrightarrow L
       \longrightarrow0.
\end{equation}
Here \(M_0\) has rank three and \(N\) has rank two.

We make the tangent identification on arbitrary Artinian tests.
For a continuous map \(A_0\to R_0\) to an Artinian local
\(k\)-algebra and a finite \(R_0\)-module \(I\), use the split
square-zero extension \(R_0\oplus I\).
Give \(I\) the divided powers \(\gamma_1(x)=x\) and
\(\gamma_j(x)=0\) for \(j\ge2\).
Grothendieck--Messing applies to this nilpotent divided-power
thickening.  The splitting provides a distinguished pullback lift;
the crystal evaluated on the thickening is the base change of
the de Rham module of that lift.
The theorem identifies deformations with lifts of the Hodge
quotient, and its full-faithfulness assertion identifies lifted
endomorphisms with operators preserving that quotient
\cite[Remark~2.3(iii), Theorem~2.4(ii),(iv)]{ChaiOort}.
Thus the strict \(\calO\)-module lifts are exactly the quotient
lifts on which \(\epsilon\) acts by zero, or equivalently lifts of
the quotient \(M_0\to L\) in \eqref{coef:hodge-quotient}.
Relative to the split lift these are parametrized by
\(\Hom(N,L\otimes I)\).
The representing deformation functor parametrizes them by
continuous derivations of \(A_0\) into \(I\).
Naturality on all these tests, and passage to the inverse limit,
give equivariantly
\[
 T_{A_0/k}\cong\Hom_{A_0}(N,L),\qquad
 \Omega^1_{A_0/k}\cong N\otimes L^{-1}.
\]
In particular
\begin{equation}\label{coef:canonical-line}
 \det\Omega^1_{A_0/k}\cong(\det M_0)L^{-3}.
\end{equation}
This is the operator-compatible version of the usual
Kodaira--Spencer calculation.

Let \(\phi\) be absolute Frobenius of \(A_0\).
Functoriality and crystalline base change give \(A_0\)-linear
maps between \(\phi^*\mathcal M\) and \(\mathcal M\) induced by
relative Frobenius and Verschiebung, with both composites multiplication
by three and hence zero
\cite[Section~2.2, Equation~(2.4), pp.~24--25]{BerthelotMessing}.
At the closed point the integral covariant Lie quotient has dimension
one, so the corresponding operator has elementary divisors five units
and one factor three.  The relation \(FV=VF=3\) gives the complementary
operator one unit and five factors three; see
\cite[Theorem~4.33(3) and Remark~4.34(2)]{ChaiOort}.
Thus one of the reduced maps, say \(U\), has rank five, and its
complement has rank one.  Their corresponding minors remain units.  Over the
fraction field of \(A_0\) the ranks are therefore at least five
and one, and their zero composition forces the first rank to
be exactly five.

Both source and target of \(U\) are free rank-three modules over
\(B_0\).  On the Frobenius twist the operator is the
\emph{pullback of the endomorphism} \(\epsilon_{\mathcal M}\).
Naturality of the crystalline maps makes \(U\) linear for these
operators: in one direction the identity is
\(U\phi^*(\epsilon_{\mathcal M})=\epsilon_{\mathcal M}U\),
and in the other direction it is the reversed identity.
These identities follow on every Artinian test from functoriality
for the endomorphism \([\pi]\), and pass to the formal base by
compatible inverse limits.  This statement does not apply absolute Frobenius
to the dual-number algebra, which would send its scalar
\(\epsilon\) to zero.

In \(B_0\)-bases write
\[
 U=U_0+\epsilon U_1,\qquad
 \det_{B_0}(U)=a+\epsilon b.
\]
Its underlying \(A_0\)-linear matrix has block form
\[
 \begin{pmatrix}U_0&0\\ U_1&U_0\end{pmatrix}.
\]
Consequently \(0=\det_{A_0}(U)=a^2\).
The ring \(A_0\) is reduced, so \(a=0\).
At the closed point, rank five forces elementary divisors
\(1,1,\epsilon\) over \(k[\epsilon]/(\epsilon^2)\).
Thus \(b\) is a unit.

This division of the determinant is intrinsic.  For a free
rank-one \(B_0\)-module, multiplication by \(\epsilon\) identifies
its reduction modulo \(\epsilon\) with its submodule multiplied
by \(\epsilon\).  The determinant map just obtained has image
in that submodule and kills \(\epsilon\) times its source.
Dividing and reducing therefore gives an isomorphism of the
reduced determinant lines.  Reversing it if necessary, we obtain
a \(\mathcal C\)-equivariant isomorphism
\[
 \phi^*(\det M_0)\xrightarrow{\ \cong\ }\det M_0.
\]
Choose a generator \(e\), with image of \(1\otimes e\) equal
to \(a_0e\), \(a_0\in A_0^\times\).
After a finite extension of \(k\), Hensel's lemma gives a unit
\(b_0\) with \(b_0^2=a_0^{-1}\).
Replacing \(e\) by \(b_0e\) makes it Frobenius-fixed.
For \(c\in\mathcal C\), write \(c(e)=\lambda_c e\).
Equivariance gives \(\lambda_c^3=\lambda_c\), and a unit
satisfying this identity in the local ring is \(1\) or \(-1\).
These scalars define a sign character \(\chi\).
The cotangent weights in \eqref{coef:canonical-line} are \(2,8\)
and the Lie weight is one.  Hence
\[
 \operatorname{wt}_T(\det M_0)=2+8+3=13\pmod{26}.
\]

It remains to relate \(M_0\) to \(V\).  Use the first
principal-parts module of \(L^{-2}\), with its exact sequence
\[
 0\longrightarrow\Omega^1_{A_0/k}\otimes L^{-2}
 \longrightarrow P^1_{A_0/k}(L^{-2})
 \longrightarrow L^{-2}\longrightarrow0.
\]
The jet map is linear over cubes, since \(d(f^3)=0\).
Thus coefficient evaluation induces an \(A_0\)-linear,
\(\mathcal C\)-equivariant map
\[
 A_0\otimes_K V\longrightarrow P^1_{A_0/k}(L^{-2}).
\]
Functoriality of principal parts includes simultaneous changes
of the base and the line trivialization, so the asserted
equivariance is independent of the chosen period.
Using the basis whose scalar evaluations are \(1,r_1,r_2\),
its matrix of values and derivatives is
\[
 \begin{pmatrix}
 1&r_1&r_2\\
 0&1&0\\
 0&0&1
 \end{pmatrix}.
\]
It is therefore an isomorphism.  Taking determinants and using
\eqref{coef:canonical-line} gives
\[
 \det(A_0\otimes_KV)
 \cong L^{-6}\det\Omega^1_{A_0/k}
 \cong L^{-9}\chi.
\]
Finally restrict to \(D'\).  The invariant element \(\theta\)
is primitive and gives a trivial subline; in principal parts
its transverse derivative is nonzero even though its value
on the curve is zero.  Taking the quotient by this subline
gives \(D'\otimes_{S_0}W\).
Its determinant is unchanged, proving
\eqref{coef:determinant-formula}.
\end{proof}

\begin{theorem}[Equivariant vanishing and invariant vectors]
\label{coef:hom-vanishing}
For every integer \(n\), including after meromorphic localization,
\[
 \Hom_{\mathcal C,S_0}(W,P_n)=0.
\]
Every nonzero invariant integral vector in \(WP_n\) evaluates as
a constant multiple of \(h^d\), where
\[
 8d=2-3n,\qquad d\ge0.
\]
If \(h\) occurs in \(WP_{-2}\), the corresponding line
\[
 \ell=hP_2\subset W
\]
is primitive and stable.  After extension to \(D'\), its quotient
has character \(w^{-15}\chi\).
More generally, if \(h^{1+3k}\) occurs in
\(WP_{-2-8k}\), \(k\ge0\), then \(h\) already occurs in
\(WP_{-2}\).
\end{theorem}

\begin{proof}
For a rank-two module,
\[
 W^\vee\cong W\otimes(\det W)^{-1}.
\]
An equivariant homomorphism \(W\to P_n\) is an invariant vector
in \(W^\vee P_n\).
Extend its determinant line to the curve and use
Lemma~\ref{coef:determinant}, together with the injective
evaluation of Theorem~\ref{coef:fixed-geometry}.
The vector then evaluates injectively in
\[
 w^{-2}w^9\chi\,w^{3n}=w^{7+3n}\chi.
\]
Tensoring the evaluation by a line over the smaller cube ring
preserves this injectivity: after extending and trivializing
the line, all values acquire the same nonzero factor.
The argument does not assert injectivity after arbitrary
extension of coefficients in the two-dimensional source.
Lemma~\ref{coef:hasse-powers} rules out the resulting invariant
section and proves the vanishing.

A vector in \(WP_n\) evaluates in \(w^{3n-2}\).
Lemma~\ref{coef:hasse-powers} gives exactly
\(8d=2-3n\), and integrality gives \(d\ge0\).
If \(h\) occurs, the generator \(q^2h\) of the indicated
submodule of \(W\) has scalar valuation one.
It cannot be divisible by a nonunit of \(S_0=k[[\tau^3]]\),
since such divisibility would raise both coefficients in
the valuation-zero/valuation-one basis by at least three.
Thus the line is primitive.
The section \(h\) is invariant and \(P_2\) extends to \(w^6\).
The line is consequently stable with extended character \(w^6\);
division of \(w^{-9}\chi\) by this character gives
\(w^{-15}\chi\) for the quotient.

For the final assertion, multiply the vector by \(q^{2+8k}\)
to regard it as an element of \(W\).  Its evaluation is
\[
 (q^2h)\,(q^8h^3)^k.
\]
Here \(q^8h^3\) is the scalar parameter \(\tau^3\) of \(S_0\).
The evaluation has order \(1+3k\).
In the basis of values of orders zero and one, distinct
residue classes modulo three prevent leading-term cancellation.
Both coefficients are therefore divisible by \(\tau^{3k}\).
Dividing within \(W\) gives the vector \(q^2h\), or
equivalently \(h\in WP_{-2}\).
\end{proof}
